% ==============================================================================
% ПАРТИЯ B03-037: DISS-005-B2012 .. DISS-005-B2046
% Глава 3: Фильтрация пустых осциллограмм (включая Таблицу 6, TAB-088)
% ==============================================================================

% DISS-005-B2012
% DISS-005-TAB-088
\paragraph{Таблица 6. Уставки для функции фильтрации «пустых» осциллограмм}\label{tab:ch03_empty_filter_settings}

\mbox{}\par\vspace{0.2cm}
{\small
\centering
\setlength{\tabcolsep}{4.5pt}
\renewcommand{\arraystretch}{1.3}
\noindent
\begin{tabular}{|>{\centering\arraybackslash}p{100pt}|>{\centering\arraybackslash}p{75pt}|>{\centering\arraybackslash}p{80pt}|>{\centering\arraybackslash}p{65pt}|>{\centering\arraybackslash}p{65pt}|}
\hline
% DISS-005-B2013
\multicolumn{2}{|c|}{сигнал} &
% DISS-005-B2014
% DISS-005-B2015
% DISS-005-B2016
\begin{tabular}[c]{@{}c@{}}Расхождение\\(delta)\end{tabular} &
% DISS-005-B2017
std\_dev &
% DISS-005-B2018
max\_abs \\ \hline
% DISS-005-B2019
\multirow{2}{*}{Нормализованные} &
% DISS-005-B2020
Ток &
% DISS-005-B2021
0.1/20 &
% DISS-005-B2022
0.05/20 &
% DISS-005-B2023
0.005/20 \\ \cline{2-5}
% DISS-005-B2024
&
% DISS-005-B2025
Напряжение &
% DISS-005-B2026
0.05/3 &
% DISS-005-B2027
0.025/3 &
% DISS-005-B2028
0.05/3 \\ \hline
% DISS-005-B2029
\multirow{2}{*}{Сырые} &
% DISS-005-B2030
Ток &
% DISS-005-B2031
0.4 &
% DISS-005-B2032
0.2 &
% DISS-005-B2033
Х \\ \cline{2-5}
% DISS-005-B2034
&
% DISS-005-B2035
Напряжение &
% DISS-005-B2036
0.05 &
% DISS-005-B2037
0.025 &
% DISS-005-B2038
Х \\ \hline
\end{tabular}
\par
}
\vspace{0.2cm}

% DISS-005-B2039
{\emergencystretch=1.5em \textit{\textbf{Примечание:} Нормализация проводится с~учётом масштабирования базисов (умножением на коэффициенты 3 для напряжений и~20 для токов, см. раздел~3.5.2), поэтому уставки для нормализованных сигналов здесь представлены как доля от этих масштабированных базисов (т.е. поделены на эти коэффициенты)}\par}

% DISS-005-B2040

% DISS-005-B2041
\noindent\textbf{\underline{Принятие решения}}\par\vspace{0.1cm}

% DISS-005-B2042
Если хотя бы один из проанализированных каналов демонстрирует активность (превышение порогов), осциллограмма считается «непустой».

% DISS-005-B2043
\noindent\textbf{\underline{Сохранение результатов}}\par\vspace{0.1cm}

% DISS-005-B2044
Список имен «непустых» осциллограмм сохраняется в указанный CSV-файл.

% DISS-005-B2045
{\emergencystretch=1.5em Применение данной функции (EmptyOscFilter) позволяет значительно сократить объем данных, требующих дальнейшей ручной или более сложной автоматической разметки, сфокусировавшись на наиболее информативных записях. Важно отметить, что пороги для этой функции требуют тщательной настройки и могут зависеть от специфики анализируемых данных и целей исследования. Использование предварительно нормализованных сигналов (если опция use\_\allowbreak norm\_\allowbreak osc = True и данные для нормализации доступны) предпочтительнее, так как позволяет применять более стабильные абсолютные пороги. Если же нормализация невозможна, анализ «сырых» данных с относительными порогами является запасным вариантом.\par}

% DISS-005-B2046
\clearpage
